\usepackage{dashrule}

\newcommand{\hmwkTitle}{Homework 5}
\newcommand{\hmwkDueDate}{Sunday, October 24, 2021}
\newcommand{\hmwkHeaderTitle}{ICCS206 Discrete Mathematic: Homework 5}
\documentclass{article}

\usepackage{fancyhdr}
\usepackage{amsmath}
\allowdisplaybreaks
\usepackage{amsthm}
\usepackage{amsfonts}
\usepackage{tikz}
\usepackage[plain]{algorithm}
\usepackage{algpseudocode}
\usepackage{enumitem}
\usepackage{graphicx}
\usepackage{hyperref}

\usetikzlibrary{automata,positioning}

%
% Basic Document Settings
%

\topmargin=-0.45in
\evensidemargin=0in
\oddsidemargin=0in
\textwidth=6.5in
\textheight=9.0in
\headsep=0.25in

\linespread{1.1}

\pagestyle{fancy}
\lhead{\hmwkAuthorName}
\chead{} % empty
\rhead{\hmwkHeaderTitle}
\cfoot{\thepage}

\renewcommand\headrulewidth{0.4pt}
\renewcommand\footrulewidth{0pt}

\setlength\parindent{0pt}

%
% Homework Details
%

\providecommand{\hmwkTitle}{Homework}
\providecommand{\hmwkDueDate}{TBD}
\providecommand{\hmwkClass}{ICCS206 Discrete Mathematic}
\providecommand{\hmwkClassTime}{Section\ 1}
\providecommand{\hmwkClassInstructor}{Pakawut Jiradilok}
\providecommand{\hmwkAuthorName}{\textbf{Jiraroj Wiruchpongsanon (6781617)}}
\providecommand{\hmwkHeaderTitle}{\hmwkClass:\ \hmwkTitle}

%
% Title Page
%

\title{
    \vspace{2in}
    \textmd{\textbf{\hmwkClass:\ \hmwkTitle}}\\
    \normalsize\vspace{0.1in}\small{Due\ on\ \hmwkDueDate}\\
    \vspace{0.1in}\large{\textit{\hmwkClassInstructor\ \hmwkClassTime}}
    \vspace{3in}
}

\author{\hmwkAuthorName}
\date{}

\renewcommand{\part}[1]{\textbf{\large Part \Alph{partCounter}}\stepcounter{partCounter}\\}

%
% Helper Commands
%

\newcounter{partCounter}
\newcommand{\solution}{\textbf{\large Solution}}

\newtheorem*{theorem}{Theorem}

\theoremstyle{remark}
\newtheorem*{remark}{Remark}

\begin{document}

\maketitle
\newpage


\section*{Problem 1}
Find closed-form expressions for the following sums/products.
\begin{enumerate}[label=\arabic*.]
    \item $\displaystyle \sum_{i=1}^n (2i + 1)$
    \item $\displaystyle \sum_{i=1}^n \sum_{j=1}^m (i + j^2)$
    \item $\displaystyle \sum_{i=1}^n \sum_{j=1}^m 2^{i+2j}$
    \item $\displaystyle \prod_{i=1}^n \prod_{j=1}^m 2^i 3^j$
\end{enumerate}

\section*{Problem 2}
Use integral bounds to give upper and lower bounds for
\begin{enumerate}[label=\arabic*.]
    \item $\displaystyle \sum_{x=1}^n \frac{1}{x^2}$
    \item $\displaystyle \sum_{x=1}^n \frac{x^2}{3}$
\end{enumerate}

\section*{Problem 3}
Bossy wants to build an $n$-story house of cards (pyramid style). Derive a formula for the number of cards required for $n$ stories, and determine how many decks are needed for a 10-story house.

\section*{Problem 4: Data Structure Cheat Sheet}
Solve or verify the following recurrences (assume $T(1)=1$).
\begin{enumerate}[label=\arabic*.]
    \item $T(n) = T(n-1) + n$ (verify by induction).
    \item $T(n) = T(n-1) + n^3$.
    \item $T(n) = T(\lfloor n/2 \rfloor) + 1$.
    \item $T(n) = T(\lfloor n/2 \rfloor) + n$.
    \item $T(n) = 2T(\lfloor n/2 \rfloor) + n$.
    \item $T(n) = 2T(\lfloor n/2 \rfloor) + 1$.
    \item (Optional) $T(n) = 2T(\lfloor n/2 \rfloor) + \log_2 n$.
\end{enumerate}

\section*{Problem 5}
For each pair $(f,g)$, determine which of $\{\sim, o, O, \omega, \Omega, \Theta\}$ apply.
\begin{enumerate}[label=\arabic*.]
    \item $f(n) = 30n + 900\log n$, $g(n) = n$
    \item $f(n) = \log n$, $g(n) = n$
    \item $f(n) = (\log n)^2$, $g(n) = n$
    \item $f(n) = \sqrt{n}$, $g(n) = (\log n)^{999}$
    \item $f(n) = n2^n$, $g(n) = n$
    \item $f(n) = n^2$, $g(n) = 1.0000000001^n$
    \item $f(n) = 200{,}000$, $g(n) = 1$
    \item $f(n) = 2^n$, $g(n) = 10^n$
\end{enumerate}

\section*{Problem 6}
Solve the recurrences:
\begin{enumerate}[label=\arabic*.]
    \item $T_{i+1} = 5T_i - 6T_{i-1}$ with $T_0 = 8$, $T_1 = 17$
    \item $T_{i+1} = -T_i + 12T_{i-1}$ with $T_0 = 0$, $T_1 = -7$
\end{enumerate}

\medskip
\solution

\begin{enumerate}
    \item we'd want to write the $i^{th}$ term of $T_i$ in a form of $x^i$, so we let

        \[T_i = x^i\]

        thus 

        \[T_{i+1} = 5T_i - 6T_{i-1} \equiv x^{i+1} = 5x^i - 6x^{i-1}\]

        divide the newly created equation with x^{i-1} gives

        \[x^2 - 5x + 6 = 0\]

        and we solve this characteristic equation with quadratic formula

        \begin{align*}
            x &= \frac{-(-5) \pm \sqrt{(-5)^2 - 4(1)(6))}}{2(1)} \\
            x &= \frac{5 + 1}{2}, \frac{5 - 1}{2} \\
            x &= 2, 3
        \end{align*}

        the solution could be formed from $2^i$ and $3^i$  by adjusting the two constant attached to them

        \[T_i = c_1 2^i + c_2 3^i\]

        we find those two constant from the given $T_0$ and $T_1$ as

        \begin{align*}
        T_0 &= 8 = c_1 2^0 + c_2 3^0 \\
            &= 8 = c_1 + c_2 \tag{i} 
        \end{align*}

        and 

        \begin{align*}
        T_1 &= 17 = c_1 2^1 + c_2 3^1 \\
            &= 17 = 2c_1 + 3c_2 \tag{ii}
        \end{align*}

        solving this linear system gives $c_1 = 7 \text{, } c_2 = 1$, and our closed form solution is complete

        \[\boxed{T_{i+1} = 7 \cdot 2^{i+1} + 3^{i+1}}\]

        \[\cong \boxed{T_{i} = 7 \cdot 2^{i} + 3^{i}}\]

    \medskip
    \begin{center}
    \rule{0.3\linewidth}{0.2pt}
    \end{center}
    \medskip    

    \item we'd want to write the $i^{th}$ term of $T_i$ in a form of $x^i$, so we let

        \[T_i = x^i\]

        thus 

        \[T_{i+1} = -T_i + 12T_{i-1} \equiv x^{i+1} = -x^i + 12x^{i-1}\]

        divide the newly created equation with x^{i-1} gives

        \[x^2 + x - 12 = 0\]

        and we solve this characteristic equation with quadratic formula

        \begin{align*}
            x &= \frac{-1 \pm \sqrt{(1)^2 - 4(1)(-12)}}{2(1)} \\
            x &= \frac{-1 + 7}{2}, \frac{-1 - 7}{2} \\
            x &= 3, -4
        \end{align*}

        the solution could be formed from $3^i$ and $(-4)^i$  by adjusting the two constant attached to them

        \[T_i = c_1 3^i + c_2 (-4)^i\]

        we find those two constant from the given $T_0$ and $T_1$ as

        \begin{align*}
        T_0 &= 0 = c_1 3^0 + c_2 (-4)^0 \\
            &= 0 = c_1 + c_2 \tag{i} 
        \end{align*}

        and 

        \begin{align*}
        T_1 &= -7 = c_1 3^1 + c_2 (-4)^1 \\
            &= -7 = 3c_1 - 4c_2 \tag{ii}
        \end{align*}

        solving this linear system gives $c_1 = -1 \text{, } c_2 = 1$, and our closed form solution is complete

        \[\boxed{T_{i+1} = -1 \cdot 3^{i+1} + (-4)^{i+1}}\]

        \[\cong \boxed{T_{i} = -1 \cdot 3^{i} + (-4)^{i}}\]
\end{enumerate}

\vspace{1em}
\hrulefill
\section*{Problem 7: Roulette Trick Redux}
Consider a deck with $r$ red cards and $b$ black cards. You always bet on red. If you win a turn you receive twice the bet. The bet-update rule is:
\begin{itemize}
    \item After a win, bet $w$ times the previous bet.
    \item After a loss, bet $l$ times the previous bet.
\end{itemize}
\begin{enumerate}[label=\arabic*.]
    \item Show that if $w + l = 2$, the result is independent of the card order.
    \item Determine the total gain/loss when $w + l = 2$ (in terms of $r,b,w,l$ and starting bet $m$).
    \item (Optional) With 20 black and 18 red cards, choose $w$ to maximize gain.
\end{enumerate}

\section*{Problem 8: Coupon Bond}
A coupon bond with face value $F$, maturity $N$ years, and annual payment $P$ pays $P$ each year, plus $F$ at maturity. If the annual interest rate is $r$, how much should you pay for the bond today?

\end{document}
